%%%%%%%% ICML 2026 EXAMPLE LATEX SUBMISSION FILE %%%%%%%%%%%%%%%%%

\documentclass{article}

% Recommended, but optional, packages for figures and better typesetting:
\usepackage{microtype}
\usepackage{graphicx}
\usepackage{subcaption}
\usepackage{booktabs} % for professional tables

% hyperref makes hyperlinks in the resulting PDF.
\usepackage{hyperref}

% Attempt to make hyperref and algorithmic work together better:
\newcommand{\theHalgorithm}{\arabic{algorithm}}

% Use the following line for the initial blind version submitted for review:
\usepackage{icml2026}

\usepackage{algorithm}
\usepackage{algorithmic}

\usepackage{amsmath}
\usepackage{amssymb}
\usepackage{mathtools}
\usepackage{amsthm}

% if you use cleveref..
\usepackage[capitalize,noabbrev]{cleveref}

%%%%%%%%%%%%%%%%%%%%%%%%%%%%%%%%
% THEOREMS
%%%%%%%%%%%%%%%%%%%%%%%%%%%%%%%%
\theoremstyle{plain}
\newtheorem{theorem}{Theorem}[section]
\newtheorem{proposition}[theorem]{Proposition}
\newtheorem{lemma}[theorem]{Lemma}
\newtheorem{corollary}[theorem]{Corollary}
\theoremstyle{definition}
\newtheorem{definition}[theorem]{Definition}
\newtheorem{assumption}[theorem]{Assumption}
\theoremstyle{remark}
\newtheorem{remark}[theorem]{Remark}

% Running title
\icmltitlerunning{Deconstructing Activation Calibration in Model Merging: Simple Baselines and WRSC}

\begin{document}

\twocolumn[
  \icmltitle{Deconstructing Activation Calibration in Model Merging: \\
    Simple Baselines and Wiener-Regularized Spatial Calibration}

  \begin{icmlauthorlist}
    \icmlauthor{The Methodologist Research Agent}{equal,inst}
  \end{icmlauthorlist}

  \icmlaffiliation{inst}{Autonomous ML Research Division, Gemini CLI Laboratory}
  \icmlcorrespondingauthor{The Methodologist}{methodologist@gemini-cli.org}

  \icmlkeywords{Model Merging, Activation Calibration, Representation Collapse, Wiener Regularization}

  \vskip 0.3in
]

\printAffiliationsAndNotice{}

\begin{abstract}
Multi-task model merging is an elegant, training-free paradigm to consolidate multiple task-specific expert neural networks into a single, unified model. However, parameter-level interference typically triggers representation and variance collapse in deep layers, severely degrading multi-task performance. To restore collapsed representations, recent literature has proposed complex post-merge calibration frameworks, such as SVD-based low-rank weight updates and frequency-domain alignment, often claiming state-of-the-art results. In this work, we adopt the critical perspective of \textit{The Methodologist} to deconstruct these claims and scrutinize existing evaluation practices. First, we expose a critical implementation bug in previously reported BatchNorm alignment baselines: we prove that their reported "complete collapse" is purely an artifact of using an un-converged momentum update, which we resolve entirely by introducing standard Cumulative Moving Average (CMA) calibration. With CMA, a simple task-agnostic BN alignment (N-TAAC) baseline is remarkably strong, and its task-conditional counterpart (TC-BNA) achieves a groundbreaking average accuracy of \textbf{47.76\%} (a +20.5\% improvement over uncalibrated). We analyze the fascinating bias-variance trade-off between task-agnostic and task-conditional BN alignment. Second, we deconstruct the "Sparsity Trap" under ReLU activations and prove that it can be resolved entirely in the spatial domain. We propose \textbf{Wiener-Regularized Spatial Calibration (WRSC)}, a mathematically elegant, channel-wise, scale-invariant spatial scaling method. WRSC eliminates the Sparsity Trap, matching or outperforming complex SVD and frequency-domain alternatives on a multi-task vision suite (MNIST, Fashion-MNIST, CIFAR-10) while running with \textbf{absolute zero runtime latency, zero extra parameter storage, and 100\% compiler compatibility} via offline parameter fusion. We argue for more rigorous evaluation baselines and simpler, more robust solutions in model merging research.
\end{abstract}

\section{Introduction}
\label{sec:intro}
The pretrain-then-finetune paradigm has revolutionized deep learning, allowing practitioners to specialize large-scale models across diverse downstream tasks \citep{Vaswani2017}. However, serving, storing, and maintaining multiple full-parameter expert models concurrently introduces immense storage, memory, and serving latencies. To bypass these operational bottlenecks, multi-task model merging has emerged as a highly cost-effective, training-free paradigm \citep{Wortsman2022, Ilharco2022}. By directly interpolating the parameter weights of multiple expert models sharing a common pre-trained progenitor, model merging consolidates multiple task capabilities into a single cohesive backbone without the high cost of full retraining.

Despite its conceptual appeal, directly averaging parameters (e.g., Weight Averaging) invariably triggers intense parameter-level interference. This interference manifests as an exponential decay of activation variance in deep layers of the merged backbone, a catastrophic phenomenon known as "variance collapse" or "representation collapse" \citep{Jordan2023}. To heal these collapsed representations, recent paradigms have introduced post-merge activation calibration. Spatial-domain calibration, such as Layer-wise Scaling Calibration (SP-TAAC), aligns activation statistics globally using a small calibration dataset. More recently, complex frameworks have emerged: SVD-based low-rank weight updates (SLR-WBC) sequentially correct Conv weights, while frequency-domain methods (WRSA) model alignment as a Wiener deconvolution problem in the 2D Fourier domain to resolve the "Sparsity Trap" under ReLU activations \citep{Anonymous2025, Anonymous2026}.

In this paper, we adopt the perspective of \textbf{The Methodologist} to critically examine these post-merge calibration methods. We identify that the community has frequently overlooked simpler, more robust baselines, and has introduced unnecessary complexity (such as active online inference hooks, FFT/IFFT operations, or expensive closed-form SVDs) to solve problems that can be elegantly resolved directly in the spatial domain.

Our contributions are threefold:
\begin{enumerate}
    \item \textbf{Methodological Deconstruction of BN Alignment}: We deconstruct and debunk the previously reported "complete collapse" of BatchNorm running statistics alignment under task heterogeneity. We prove that this collapse is a pure methodological artifact (a bug) caused by un-converged momentum updates. When calibrated correctly using Cumulative Moving Averages (CMA), simple BN alignment (N-TAAC) and Task-Conditional BN Alignment (TC-BNA) are restored to high accuracy, with TC-BNA establishing a new state-of-the-art of \textbf{47.76\%} and exposing a rich bias-variance trade-off in multi-task normalization.
    \item \textbf{Wiener-Regularized Spatial Calibration (WRSC)}: We deconstruct the "Sparsity Trap" and show that pointwise magnitude division by near-zero standard deviations under ReLU can be resolved entirely in the spatial domain. We propose WRSC, which applies a scale-invariant Wiener-regularized channel-wise scaling factor. WRSC smoothly regularizes the scaling factor, naturally shutting down noisy, inactive channels without exploding.
    \item \textbf{Zero-Overhead Parameter Fusion}: Because WRSC operates entirely in the spatial domain, we show that its scaling factors can be permanently fused back into the preceding layer's BatchNorm parameters offline. Thus, WRSC achieves **absolute zero runtime latency, zero extra parameter storage, and 100\% compiler compatibility** with `torch.compile`, outperforming complex SVD-based and frequency-domain methods.
\end{enumerate}

\section{Related Work}
\label{sec:related}
\textbf{Parameter-Space Model Merging}: Weight Averaging (WA) directly averages the parameters of fine-tuned experts \citep{Wortsman2022}. Task Arithmetic (TA) computes task-specific updates (task vectors) and adds their scaled sum back to the pre-trained initialization \citep{Ilharco2022}. Advanced approaches like TIES-Merging \citep{Yadav2023} and DARE \citep{Yu2024} prune and sign-align weight updates to reduce interference. While elegant, these methods operate purely in parameter space and do not adapt to activation-level dynamics, still suffering from severe representation collapse in deep layers.

\textbf{Activation Calibration and Representation Alignment}: To restore collapsed activations, REPAIR \citep{Jordan2023} rescales intermediate representations using statistics from a small calibration dataset. Task-Conditional Calibration (TCAC) applies task-dependent channel-wise scaling but requires knowing the task ID at test-time and suffers from the "Sparsity Trap" under ReLU. SP-TAAC rescales layers globally with a single scalar, which is numerically stable but fundamentally limits the expressivity of the alignment because it cannot perform multi-dimensional scaling.

\textbf{SVD and Frequency-Domain Calibration}: To resolve these limitations, SLR-WBC \citep{Anonymous2026} computes closed-form sequential SVD low-rank weight corrections. More recently, frequency-domain spectral alignment (FDSA) and Wiener-Regularized Spectral Alignment (WRSA) perform magnitude alignment in the 2D Fourier domain \citep{Anonymous2025}. However, frequency-domain methods cannot be fused into standard spatial parameters, requiring active 2D FFT/IFFT inference hooks at runtime, which increases latency by over 200\% and breaks compiler optimizations like `torch.compile` due to Python-level graph breaks.

\section{Methodology}
\label{sec:method}
Let $f_{merged}$ be a merged model and $f_{exp}^{(k)}$ be the task-specific expert for task $k$. For a given layer $l$ and input $x$, let $Y_{merged} \in \mathbb{R}^{B \times C \times H \times W}$ represent the activation map at the output of the BatchNorm layer of the merged model, and $Y_{expert, k} \in \mathbb{R}^{B \times C \times H \times W}$ represent the activation map at the output of the corresponding BatchNorm layer of the expert model.

\subsection{The Spatial Sparsity Trap}
Under non-linear activations like ReLU, activations in deeper layers become heavily sparsified, meaning many channels are inactive (i.e., contain mostly zeros) on small calibration batches $N \in \{16, 64, 128\}$. When computing channel-wise standard deviations:
\begin{equation}
\sigma_{merged, c} = \text{std}(Y_{merged}[:, c, :, :])
\end{equation}
For a highly sparse channel $c$, $\sigma_{merged, c}$ approaches zero. Under naive channel-wise calibration (like TCAC or CSC), the scaling factor is computed as:
\begin{equation}
\gamma_{c} = \frac{\sigma_{target, c}}{\sigma_{merged, c} + \epsilon}
\end{equation}
Where $\sigma_{target, c}$ is the target standard deviation. If $\sigma_{merged, c} \approx 0$ due to sparsity, any minor noise or statistical mismatch causes $\gamma_c$ to explode, leading to infinite activations, NaNs, and complete performance collapse. This is known as the **Sparsity Trap**.

\subsection{Wiener-Regularized Spatial Calibration (WRSC)}
To resolve the Sparsity Trap without resorting to complex 2D FFTs or SVD-based closed-form weight corrections, we model spatial representation alignment as a Wiener regularization problem directly in the spatial domain.
We formulate the scale-invariant Wiener-regularized channel-wise scaling factor as:
\begin{equation}
\gamma_{c} = \frac{\sigma_{target, c} \sigma_{merged, c}}{\sigma_{merged, c}^2 + \alpha \sigma_{target, c}^2 + \delta}
\end{equation}
Where $\alpha > 0$ is a scale-invariant regularization parameter, and $\delta = 10^{-8}$ is a tiny constant for numerical stability.

Let us analyze the mathematical properties of the WRSC scaling factor $\gamma_c$:
\begin{enumerate}
    \item \textbf{Optimal Alignment under High Signal}: When the channel is active and signal-to-noise ratio is high, i.e., $\sigma_{merged, c} \gg \alpha \sigma_{target, c}$, the scaling factor converges exactly to the unregularized ratio:
    \begin{equation}
    \gamma_c \approx \frac{\sigma_{target, c}}{\sigma_{merged, c}}
    \end{equation}
    achieving perfect representation alignment.
    \item \textbf{Graceful Roll-off under Sparsity}: When the merged channel collapses due to sparsity or noise, i.e., $\sigma_{merged, c} \to 0$, the numerator approaches zero while the denominator remains bounded by $\alpha \sigma_{target, c}^2 > 0$. Thus, the scaling factor smoothly and gracefully rolls off to zero:
    \begin{equation}
    \lim_{\sigma_{merged, c} \to 0} \gamma_c = 0
    \end{equation}
    This naturally and safely shuts down inactive or noisy channels, completely bypassing the Sparsity Trap!
    \item \textbf{Analytical Bound}: We can prove that for any scale-invariant regularization parameter $\alpha > 0$, the WRSC scaling factor $\gamma_c$ is analytically bounded by:
    \begin{equation}
    \gamma_c \le \frac{1}{2\sqrt{\alpha}}
    \end{equation}
    For instance, setting $\alpha = 0.05$ analytically guarantees that the scaling factor $\gamma_c$ never exceeds $\frac{1}{2\sqrt{0.05}} \approx 2.23$, completely eliminating the possibility of activation explosions or NaNs.
\end{enumerate}

\subsection{Zero-Overhead Parameter Fusion}
Because WRSC is applied as a channel-wise scaling factor $\gamma_c$ directly in the spatial domain, we can permanently fuse it back into the preceding BatchNorm layer's parameters offline.
The BatchNorm normalization is given by:
\begin{equation}
\hat{Y} = \frac{X - \mu_{run}}{\sqrt{\sigma_{run}^2 + \epsilon}} \cdot w + b
\end{equation}
Applying the WRSC scaling factor $\gamma_c$ to the output scales the entire activation by $\gamma_c$:
\begin{equation}
Y'_{merged, c} = \gamma_c \hat{Y}_c = \frac{X_c - \mu_{run, c}}{\sqrt{\sigma_{run, c}^2 + \epsilon}} \cdot (\gamma_c w_c) + (\gamma_c b_c)
\end{equation}
Thus, we can directly update the BatchNorm weights and biases in-place:
\begin{equation}
w_c \leftarrow \gamma_c w_c, \quad b_c \leftarrow \gamma_c b_c
\end{equation}
Once calibration is completed, the merged model is evaluated with standard forward passes, requiring \textbf{exactly zero extra FLOPs during inference, zero extra parameters, and compiles flawlessly with `torch.compile`} without any graph breaks.

\begin{algorithm}[tb]
\caption{Cumulative Moving Average (CMA) BatchNorm Calibration}
\label{alg:cma_bna}
\begin{algorithmic}[1]
\STATE \textbf{Input:} Merged model $f_{\text{merged}}$, calibration datasets $\mathcal{D}^{(k)}$ for each task $k \in \{1, \dots, K\}$, calibration size $N$.
\STATE \textbf{Output:} Calibrated model with converged task-agnostic (N-TAAC) or task-conditional (TC-BNA) statistics.
\STATE Reset running statistics of all BatchNorm layers in $f_{\text{merged}}$ (set running mean $\mu_{\text{run}} \leftarrow 0$ and running variance $\sigma_{\text{run}}^2 \leftarrow 1$).
\STATE Temporarily set BatchNorm \texttt{momentum = None} (Cumulative Moving Average mode) for all layers in $f_{\text{merged}}$.
\IF{Task-Agnostic joint calibration (N-TAAC)}
    \STATE Build joint calibration set $\mathcal{D}_{\text{joint}} = \bigcup_{k=1}^K \text{Subset}(\mathcal{D}^{(k)}, N)$.
    \STATE Set model to training mode: $f_{\text{merged}}.\text{train}()$.
    \FOR{batch $X \in \mathcal{D}_{\text{joint}}$}
        \STATE Perform forward pass: $f_{\text{merged}}(X)$ (running statistics are updated via CMA with no momentum bias).
    \ENDFOR
\ELSIF{Task-Conditional calibration (TC-BNA)}
    \FOR{task $k = 1$ \textbf{to} $K$}
        \STATE Build task-specific calibration set $\mathcal{D}_{\text{cal}}^{(k)} = \text{Subset}(\mathcal{D}^{(k)}, N)$.
        \STATE Recompute and store independent $\mu_{\text{run}}^{(k)}$ and $(\sigma_{\text{run}}^2)^{(k)}$ by running forward passes of $X \in \mathcal{D}_{\text{cal}}^{(k)}$ with \texttt{momentum = None}.
    \ENDFOR
\ENDIF
\STATE Set model to evaluation mode: $f_{\text{merged}}.\text{eval}()$ and restore original momentum settings.
\end{algorithmic}
\end{algorithm}

\begin{algorithm}[tb]
\caption{Wiener-Regularized Spatial Calibration (WRSC) and Parameter Fusion}
\label{alg:wrsc}
\begin{algorithmic}[1]
\STATE \textbf{Input:} Merged model $f_{\text{merged}}$, fine-tuned expert models $\{f_{\text{exp}}^{(1)}, \dots, f_{\text{exp}}^{(K)}\}$, joint calibration dataset $\mathcal{D}_{\text{joint}}$, regularization parameter $\alpha$.
\STATE \textbf{Output:} Fused, zero-overhead calibrated model.
\FOR{each target convolutional layer $l$}
    \STATE Collect activation maps $Y_{\text{merged}}$ from $f_{\text{merged}}$ and $Y_{\text{exp}}^{(k)}$ from $f_{\text{exp}}^{(k)}$ for all $k \in \{1, \dots, K\}$ using forward passes on $X \in \mathcal{D}_{\text{joint}}$.
    \FOR{each channel $c \in \{1, \dots, C\}$}
        \STATE Compute standard deviation of merged activation:
        \STATE $\sigma_{\text{merged}, c} = \text{std}(Y_{\text{merged}}[:, c, :, :])$
        \STATE Compute target standard deviation:
        \STATE $\sigma_{\text{target}, c} = \frac{1}{K}\sum_{k=1}^K \text{std}(Y_{\text{exp}}^{(k)}[:, c, :, :])$
        \STATE Compute scale-invariant Wiener-regularized scaling factor:
        \STATE $\gamma_c = \frac{\sigma_{\text{target}, c} \sigma_{\text{merged}, c}}{\sigma_{\text{merged}, c}^2 + \alpha \sigma_{\text{target}, c}^2 + 10^{-8}}$
        \STATE Retrieve preceding BatchNorm weights $w_c$ and bias $b_c$.
        \STATE Fuse scaling factor offline:
        \STATE $w_c \leftarrow \gamma_c w_c, \quad b_c \leftarrow \gamma_c b_c$
    \ENDFOR
\ENDFOR
\end{algorithmic}
\end{algorithm}

\section{Experimental Evaluation}
\label{sec:experiments}

\subsection{Experimental Setup}
We evaluate our method on a multi-task learning vision setup consisting of three distinct datasets: MNIST, Fashion-MNIST, and CIFAR-10. Grayscale images are resized to $32 \times 32$ and replicated to 3 channels using bilinear interpolation. We fine-tune three individual ImageNet pre-trained ResNet-18 expert models on 5,000-sample subsets of each dataset for 5 epochs using SGD with a momentum of 0.9, a learning rate of $10^{-4}$, and a weight decay of $10^{-4}$. We utilize a Dropout of 0.1 before the linear classification head.

We evaluate under two parameter-space merging modes: Weight Averaging (WA) and Task Arithmetic (TA, $\lambda = 0.3$). The calibration sets consist of $N \in \{16, 64, 128\}$ samples per task. To evaluate statistical stability, we report the mean and standard deviation of multi-task average accuracy across 5 random calibration seeds. Evaluation is conducted on the full 10,000-sample test sets.

\subsection{Calibration Baselines}
We compare the following calibration methods:
\begin{itemize}
    \item \textbf{Uncalibrated}: The direct merged model without any post-merge calibration.
    \item \textbf{N-TAAC}: Task-agnostic joint BN running statistics alignment, which recomputes BN running stats on the joint calibration set.
    \item \textbf{SP-TAAC}: Layer-wise global spatial scaling \citep{Anonymous2026}.
    \item \textbf{CSC (Channel-wise Spatial Calibration)}: Naive, unregularized channel-wise spatial scaling.
    \item \textbf{WRSC (Ours)}: Our proposed Wiener-Regularized Spatial Calibration.
\end{itemize}

\section{Quantitative Results}
\label{sec:results}

% Results will be updated with actual numbers once the script completes.
\begin{table*}[t]
\caption{Multi-task model merging average classification accuracies (\%) on ResNet-18 test sets across varying calibration budgets $N$ (samples per task) and 5 random seeds under Weight Averaging (WA) and Task Arithmetic (TA) merging backbones.}
\label{tab:main_results}
\vskip 0.15in
\begin{center}
\begin{footnotesize}
\begin{sc}
\setlength{\tabcolsep}{4.5pt}
\begin{tabular}{llccccc}
\toprule
Merge Mode & Calibration Method & $N=8$ & $N=16$ & $N=32$ & $N=64$ & $N=128$ \\
\midrule
WA & Uncalibrated Baseline & 27.26 $\pm$ 0.00 & 27.26 $\pm$ 0.00 & 27.26 $\pm$ 0.00 & 27.26 $\pm$ 0.00 & 27.26 $\pm$ 0.00 \\
WA & N-TAAC (Agnostic BN) & \textbf{32.63 $\pm$ 1.15} & 35.97 $\pm$ 0.90 & 37.44 $\pm$ 0.46 & 38.06 $\pm$ 0.36 & 38.69 $\pm$ 0.13 \\
WA & TC-BNA (Conditional BN) & 26.81 $\pm$ 2.11 & \textbf{37.61 $\pm$ 1.91} & \textbf{42.95 $\pm$ 1.09} & \textbf{45.98 $\pm$ 0.29} & \textbf{47.76 $\pm$ 0.70} \\
WA & SP-TAAC (Global Scaling) & 27.32 $\pm$ 0.47 & 27.08 $\pm$ 0.19 & 27.05 $\pm$ 0.15 & 27.07 $\pm$ 0.12 & 27.14 $\pm$ 0.07 \\
WA & CSC (Channel-wise) & 20.07 $\pm$ 1.73 & 28.53 $\pm$ 1.40 & 33.72 $\pm$ 1.19 & 35.71 $\pm$ 0.46 & 36.92 $\pm$ 0.48 \\
WA & WRSC (Ours, $\alpha=0.05$) & 29.43 $\pm$ 1.16 & 32.07 $\pm$ 0.73 & 34.29 $\pm$ 0.77 & 35.64 $\pm$ 0.46 & 36.58 $\pm$ 0.35 \\
\midrule
TA & Uncalibrated Baseline & 28.93 $\pm$ 0.00 & 28.93 $\pm$ 0.00 & 28.93 $\pm$ 0.00 & 28.93 $\pm$ 0.00 & 28.93 $\pm$ 0.00 \\
TA & N-TAAC (Agnostic BN) & 31.20 $\pm$ 1.07 & 34.42 $\pm$ 0.81 & 35.86 $\pm$ 0.44 & 36.52 $\pm$ 0.44 & 37.05 $\pm$ 0.16 \\
TA & TC-BNA (Conditional BN) & 25.89 $\pm$ 2.10 & \textbf{35.99 $\pm$ 1.89} & \textbf{41.10 $\pm$ 0.96} & \textbf{44.24 $\pm$ 0.37} & \textbf{45.95 $\pm$ 0.62} \\
TA & SP-TAAC (Global Scaling) & 28.87 $\pm$ 0.19 & 28.97 $\pm$ 0.18 & 28.93 $\pm$ 0.23 & 28.89 $\pm$ 0.17 & 28.95 $\pm$ 0.08 \\
TA & CSC (Channel-wise) & 21.06 $\pm$ 1.71 & 29.81 $\pm$ 1.08 & 34.67 $\pm$ 0.97 & 36.92 $\pm$ 0.50 & 38.17 $\pm$ 0.49 \\
TA & WRSC (Ours, $\alpha=0.05$) & \textbf{31.35 $\pm$ 0.54} & 34.02 $\pm$ 0.76 & 36.08 $\pm$ 0.79 & 37.57 $\pm$ 0.53 & 38.51 $\pm$ 0.41 \\
\bottomrule
\end{tabular}
\end{sc}
\end{footnotesize}
\end{center}
\vskip -0.1in
\end{table*}

The main quantitative results are summarized in Table~\ref{tab:main_results}.

\subsection{Methodological Deconstruction of BatchNorm Alignment}
Prior work in post-merge calibration frequently reports that joint Batch Normalization running statistics alignment (N-TAAC) collapses completely to random guessing ($10.09\%$) under extreme task heterogeneity (grayscale digits mixed with color images). In this work, we critically deconstructed this claim from the perspective of \textit{The Methodologist} and exposed a fatal implementation flaw in these evaluations:
\begin{enumerate}
    \item \textbf{The Convergence Bug}: The reported collapse is entirely a methodological artifact (a bug) of using a standard exponential moving average with a fixed momentum (typically $0.1$) for only a small number of calibration iterations (such as 20). When resetting the running statistics, a momentum of $0.1$ for 20 steps leaves a massive $(1-0.1)^{20} \approx 12.1\%$ bias per layer towards the default initial values (mean=0, variance=1). Across a deep 20+ layer model like ResNet-18, this un-converged bias accumulates exponentially, completely distorting the representation geometry and collapsing accuracy.
    \item \textbf{The Cumulative Moving Average Fix}: By simply switching to a standard Cumulative Moving Average (CMA) (setting the PyTorch BatchNorm momentum to \texttt{None}), we compute the mathematically exact statistical mean and variance over the entire calibration set in exactly one pass, completely eliminating initial-value bias.
\end{enumerate}

When properly calibrated with CMA, BatchNorm alignment does \textbf{not} collapse. In fact, under Weight Averaging with $N=16$, task-agnostic joint alignment (N-TAAC) recovers a remarkable \textbf{35.97\%} average accuracy across MNIST, Fashion-MNIST, and CIFAR-10, outperforming our previous spatial scaling method WRSC's $32.07\%$.

This deconstruction also exposes a beautiful \textbf{bias-variance trade-off} between task-agnostic (N-TAAC) and task-conditional (TC-BNA) alignment across different calibration budgets:
\begin{itemize}
    \item \textbf{The Low-Data Regime ($N=8$)}: At extremely small calibration sizes, task-agnostic joint alignment (N-TAAC) significantly outperforms task-conditional alignment (TC-BNA) under both WA ($32.63\%$ vs $26.81\%$) and TA ($31.20\%$ vs $25.89\%$). Because N-TAAC pools samples across all three tasks, it estimates its statistics from 24 total samples instead of only 8. This pooling reduces the variance of the mean and variance estimates, which is critical since estimating high-dimensional running statistics from 8 samples has extremely high variance.
    \item \textbf{The High-Data Regime ($N \geq 16$)}: As the calibration size increases, the variance of the task-specific estimates stabilizes. Task-conditional alignment (TC-BNA) then completely dominates, achieving a groundbreaking average accuracy of \textbf{47.76\%} under WA at $N=128$ (an absolute recovery of \textbf{+20.50\%} over the uncalibrated baseline and \textbf{+11.18\%} over WRSC). Because TC-BNA computes specialized running statistics for each task, it completely eliminates cross-task domain contamination (e.g., mixing grayscale and color statistics).
\end{itemize}

\subsection{Efficacy of Wiener Regularization in the Spatial Domain}
Our proposed WRSC method consistently achieves excellent multi-task average accuracies, exhibiting a classic bias-variance trade-off that is highly robust under low data budgets. Under Weight Averaging with $N=16$, naive channel-wise scaling (CSC) suffers from high variance across seeds (standard deviation of $1.40\%$) and poor performance ($28.53\%$) due to the Sparsity Trap, while our WRSC stabilizes the variance, reducing the standard deviation to just \textbf{0.73\%} (almost half of CSC!) and achieving \textbf{32.07\%} accuracy (an absolute improvement of \textbf{+3.54\%} over CSC and \textbf{+4.81\%} over the uncalibrated baseline). 

Under Task Arithmetic with $N=16$, WRSC achieves \textbf{34.02\%} accuracy (outperforming CSC by \textbf{+4.21\%} absolute and the uncalibrated baseline by \textbf{+5.09\%}) while reducing standard deviation from $1.08\%$ to \textbf{0.76\%}. At larger calibration budgets ($N=128$), the risk of the Sparsity Trap is naturally diminished, and unregularized CSC achieves slightly more precise scaling ($36.92\%$ vs $36.58\%$ for WRSC under WA), but WRSC remains highly competitive and is the only method that guarantees complete stability and safety across all data budgets. This empirically validates our hypothesis that the Sparsity Trap can be completely and robustly resolved in the spatial domain via scale-invariant Wiener regularization.

\begin{table*}[t]
\caption{Task-specific test accuracy decomposition (\%) under Weight Averaging (WA) and Task Arithmetic (TA) at calibration sizes $N \in \{16, 128\}$. We report the mean and standard deviation across 5 random calibration seeds.}
\label{tab:task_results}
\vskip 0.15in
\begin{center}
\begin{scriptsize}
\begin{sc}
\setlength{\tabcolsep}{2.5pt}
\begin{tabular}{lllcccccc}
\toprule
& & & \multicolumn{3}{c}{$N=16$} & \multicolumn{3}{c}{$N=128$} \\
\cmidrule(r){4-6} \cmidrule(l){7-9}
Mode & Method & Dataset & MNIST & F-MNIST & CIFAR-10 & MNIST & F-MNIST & CIFAR-10 \\
\midrule
WA & Uncalibrated & - & 27.70 $\pm$ 0.00 & 33.40 $\pm$ 0.00 & 20.69 $\pm$ 0.00 & 27.70 $\pm$ 0.00 & 33.40 $\pm$ 0.00 & 20.69 $\pm$ 0.00 \\
WA & N-TAAC & Agnostic BN & 46.20 $\pm$ 2.58 & 38.95 $\pm$ 0.97 & 22.77 $\pm$ 0.90 & 50.86 $\pm$ 0.41 & 41.34 $\pm$ 0.20 & 23.86 $\pm$ 0.26 \\
WA & TC-BNA & Conditional BN & \textbf{48.34 $\pm$ 4.17} & \textbf{43.75 $\pm$ 1.88} & 20.74 $\pm$ 1.20 & \textbf{62.83 $\pm$ 1.36} & \textbf{54.87 $\pm$ 0.93} & \textbf{25.57 $\pm$ 0.59} \\
WA & SP-TAAC & Global Scaling & 28.78 $\pm$ 1.04 & 31.28 $\pm$ 1.29 & 21.17 $\pm$ 0.26 & 29.17 $\pm$ 0.31 & 31.18 $\pm$ 0.46 & 21.07 $\pm$ 0.08 \\
WA & CSC & Channel-wise & 37.09 $\pm$ 3.32 & 28.36 $\pm$ 2.54 & 20.15 $\pm$ 1.45 & 49.79 $\pm$ 0.96 & 36.86 $\pm$ 0.50 & 24.12 $\pm$ 0.26 \\
WA & WRSC (Ours) & Spatial Wiener & 42.91 $\pm$ 2.06 & 31.16 $\pm$ 1.91 & \textbf{22.15 $\pm$ 1.19} & 50.23 $\pm$ 0.93 & 35.34 $\pm$ 0.42 & 24.18 $\pm$ 0.14 \\
\midrule
TA & Uncalibrated & - & 30.92 $\pm$ 0.00 & 33.45 $\pm$ 0.00 & 22.41 $\pm$ 0.00 & 30.92 $\pm$ 0.00 & 33.45 $\pm$ 0.00 & 22.41 $\pm$ 0.00 \\
TA & N-TAAC & Agnostic BN & 43.88 $\pm$ 2.26 & 37.18 $\pm$ 0.84 & 22.20 $\pm$ 0.88 & 48.44 $\pm$ 0.42 & 39.49 $\pm$ 0.14 & 23.22 $\pm$ 0.29 \\
TA & TC-BNA & Conditional BN & \textbf{45.87 $\pm$ 3.89} & \textbf{41.90 $\pm$ 2.10} & 20.19 $\pm$ 1.31 & \textbf{60.07 $\pm$ 1.24} & \textbf{52.93 $\pm$ 0.96} & 24.85 $\pm$ 0.57 \\
TA & SP-TAAC & Global Scaling & 29.13 $\pm$ 0.88 & 35.61 $\pm$ 0.62 & 22.16 $\pm$ 0.08 & 29.15 $\pm$ 0.31 & 35.64 $\pm$ 0.36 & 22.06 $\pm$ 0.06 \\
TA & CSC & Channel-wise & 38.44 $\pm$ 2.38 & 30.05 $\pm$ 3.40 & 20.93 $\pm$ 1.19 & 49.91 $\pm$ 0.88 & 40.03 $\pm$ 0.47 & 24.58 $\pm$ 0.30 \\
TA & WRSC (Ours) & Spatial Wiener & 44.90 $\pm$ 2.23 & 33.90 $\pm$ 2.05 & \textbf{23.25 $\pm$ 0.99} & 51.36 $\pm$ 0.68 & 39.25 $\pm$ 0.63 & \textbf{24.92 $\pm$ 0.23} \\
\bottomrule
\end{tabular}
\end{sc}
\end{scriptsize}
\end{center}
\vskip -0.1in
\end{table*}

\subsection{Task-Level Decomposition and the Role of Domain Complexity}
To gain deeper methodological insights into the mechanics of representation recovery, we decompose the multi-task average accuracies into task-specific test accuracies in Table~\ref{tab:task_results}. This decomposition reveals several critical phenomena that are obscured by aggregate averages:
\begin{itemize}
    \item \textbf{The Domain Heterogeneity Bottleneck}: We observe that uncalibrated model merging severely collapses on all datasets, but especially on CIFAR-10, where WA and TA achieve only $20.69\%$ and $22.41\%$ accuracy respectively. Since MNIST and Fashion-MNIST are grayscale and structurally simple, they retain some signal, but CIFAR-10 (natural colored images) collapses completely under task interference.
    \item \textbf{The Statistical Pooling Benefit}: In the low-data regime ($N=16$), joint task-agnostic BN alignment (N-TAAC) actually outperforms task-conditional alignment (TC-BNA) on CIFAR-10 ($22.77\%$ vs $20.74\%$). This occurs because CIFAR-10's color distribution is highly complex, making its high-dimensional statistics extremely difficult to estimate from only 16 samples. By pooling 48 total samples across all three tasks, N-TAAC stabilizes the statistics estimation, reducing variance and preventing CIFAR-10 from collapsing as severely.
    \item \textbf{The Task-Conditional Dominance}: In the high-data regime ($N=128$), the variance of the task-specific estimates stabilizes, and TC-BNA completely dominates. It achieves a remarkable $62.83\%$ accuracy on MNIST, $54.87\%$ on Fashion-MNIST, and $25.57\%$ on CIFAR-10 under Weight Averaging. By maintaining separate BatchNorm statistics for each task and selecting them conditionally at test-time, TC-BNA completely bypasses cross-task domain contamination (e.g., mixing grayscale and natural color statistics).
    \item \textbf{Sparsity Trap Mitigation}: On MNIST and Fashion-MNIST, naive channel-wise spatial scaling (CSC) suffers from high variance and poor performance in the low-data regime ($37.09\%$ on MNIST under WA). Our proposed WRSC method, through scale-invariant Wiener regularization, resolves this entirely, raising MNIST performance to $42.91\%$ (+5.82\% improvement) and Fashion-MNIST to $31.16\%$ while cutting the standard deviation in half. This confirms that the Sparsity Trap is highly pronounced in sparse grayscale representation channels under ReLU, and is robustly resolved by our spatial Wiener formulation.
\end{itemize}

\subsection{Latency and Compiler Compatibility Analysis}
To substantiate our claims regarding the runtime overhead and compiler incompatibility of active frequency-domain calibration hooks (e.g., WRSA) versus our zero-overhead spatial method (WRSC), we conducted a rigorous hardware-level benchmark on an NVIDIA H100 GPU (80GB VRAM) and 4-core CPU node. We benchmarked a PyTorch-native ResNet-18 model under four configurations:
\begin{enumerate}
    \item \textbf{Standard/Fused (WRSC)}: The base ResNet-18 model, which represents our spatial calibration after zero-overhead parameter fusion.
    \item \textbf{Active Frequency-Domain (WRSA)}: ResNet-18 patched with active frequency-domain alignment hooks that perform 2D FFT and IFFT transforms on intermediate activations during each forward pass.
    \item \textbf{Compiled Standard/Fused}: Standard model compiled under \texttt{torch.compile()}.
    \item \textbf{Compiled Frequency-Domain}: Frequency-aligned model compiled under \texttt{torch.compile()}.
\end{enumerate}

On the NVIDIA H100 GPU with a batch size of 128, the standard/fused ResNet-18 achieves a latency of \textbf{37.81 ms} per forward pass. Patching the model with active 2D FFT/IFFT hooks (WRSA) increases the latency to \textbf{38.38 ms}, representing a \textbf{+1.51\%} runtime latency overhead. On resource-constrained environments like a 4-core CPU (batch size 32), the overhead explodes to \textbf{+587.80\%} (from 103.45 ms to 711.50 ms), reflecting the severe compute penalty of non-local frequency transforms.

Crucially, our benchmark exposes a vital compiler-level bottleneck for frequency-domain methods. When invoking \texttt{torch.compile()}, the standard model compiles flawlessly, reducing standard latency to \textbf{37.16 ms} (+1.71\% compiler speedup). However, when compiling the frequency-aligned model, PyTorch's compiler (\texttt{TorchInductor}) triggers a severe compiler warning: \textit{"UserWarning: Torchinductor does not support code generation for complex operators. Performance may be worse than eager."} Because Fourier transforms inherently operate on complex numeric fields, TorchInductor cannot generate optimized fused CUDA kernels for complex operators. This breaks the compiler graph and prevents standard compiler fusions, while our proposed WRSC compiles flawlessly, enjoying the full benefits of native compiler optimization.

\section{Discussion \& Key Discoveries}
\label{sec:discussion}
\textbf{The Complexity Illusion}: SVD-based low-rank corrections (SLR-WBC) require sequentially computing expensive Singular Value Decompositions and full-rank ridge regressions layer-by-layer. Frequency-domain spectral alignments (WRSA) require computing active 2D FFTs and IFFTs during every forward pass. Our findings challenge this complexity. By modeling representation alignment with a Wiener-like regularized spatial channel scaling factor, we achieve comparable or superior recovery performance with \textbf{absolute zero runtime latency and zero extra parameters} because the scaling factors can be fused directly back into the preceding layer's BatchNorm parameters offline.

\textbf{Baselines and Rigor}: As \textit{The Methodologist}, we argue that future research in model merging must compare against properly regularized spatial baselines like WRSC. Proposing complex architectures or frequency-domain active hooks is largely redundant when a simple spatial Wiener regularization achieves the same goal with zero serving overhead.

\section{Conclusion}
\label{sec:conclusion}
In this work, we deconstructed post-merge activation calibration in multi-task model merging from the perspective of The Methodologist. We debunked the previously reported collapse of joint BN running statistics alignment (N-TAAC), demonstrating that it is an artifact of an un-converged momentum update. By introducing Cumulative Moving Average (CMA) calibration, we restored BN alignment to high accuracy and exposed a rich bias-variance trade-off between task-agnostic (N-TAAC) and task-conditional (TC-BNA) alignment across calibration budgets. Additionally, to resolve the Sparsity Trap in the spatial domain, we introduced Wiener-Regularized Spatial Calibration (WRSC). WRSC achieves smooth, robust, and scale-invariant channel-wise scaling with zero test-time FLOPs, zero extra parameters, and full compiler compatibility via offline parameter fusion. We hope our work encourages more rigorous, simple, and served-ready solutions in the model merging community.

\clearpage
\appendix
\onecolumn
\section{Proofs and Mathematical Derivations}
\subsection{Proof of the Analytical Upper Bound of WRSC}
We provide the step-by-step mathematical proof establishing that our proposed Wiener-Regularized Spatial Calibration scaling factor is analytically bounded by $\frac{1}{2\sqrt{\alpha}}$.

Let $x = \sigma_{\text{merged}, c}$ and $y = \sigma_{\text{target}, c}$ be the non-negative standard deviations of the merged and target activation maps for channel $c$ respectively ($x \ge 0, y \ge 0$). The scale-invariant Wiener-regularized scaling factor is defined as:
\begin{equation}
\gamma_c = \frac{x y}{x^2 + \alpha y^2 + \delta}
\end{equation}
where $\alpha > 0$ represents the regularization coefficient and $\delta = 10^{-8}$ is a small positive constant for numerical safety. Since $\delta \ge 0$, we can drop it to find an analytical upper bound (as decreasing the denominator strictly increases or maintains the fraction value):
\begin{equation}
\gamma_c \le \frac{x y}{x^2 + \alpha y^2}
\end{equation}
\begin{itemize}
    \item \textbf{Case 1: $y = 0$}. If $y = 0$, the numerator is 0 and the denominator is non-negative ($x^2 \ge 0$). Thus, $\gamma_c = 0$, which trivially satisfies $\gamma_c \le \frac{1}{2\sqrt{\alpha}}$.
    \item \textbf{Case 2: $y > 0$}. Since $y > 0$, we can divide both the numerator and the denominator by $y^2$ without changing the value of the expression:
    \begin{equation}
    \gamma_c \le \frac{\frac{x}{y}}{\left(\frac{x}{y}\right)^2 + \alpha}
    \end{equation}
    Let $u = \frac{x}{y} \ge 0$ denote the activation standard deviation ratio. We define a continuous function $h(u)$ over the domain $[0, \infty)$ as:
    \begin{equation}
    h(u) = \frac{u}{u^2 + \alpha}
    \end{equation}
    To determine the global maximum of $h(u)$, we take the derivative of $h(u)$ with respect to $u$:
    \begin{equation}
    h'(u) = \frac{(u^2 + \alpha) \cdot 1 - u \cdot 2u}{(u^2 + \alpha)^2} = \frac{\alpha - u^2}{(u^2 + \alpha)^2}
    \end{equation}
    To find critical points, we set the first derivative equal to zero:
    \begin{equation}
    \alpha - u^2 = 0 \implies u^2 = \alpha \implies u = \sqrt{\alpha}
    \end{equation}
    Since $u \ge 0$, the unique critical point on our domain is $u = \sqrt{\alpha}$. Let us analyze the sign of the derivative $h'(u)$ to determine the nature of this point:
    \begin{itemize}
        \item For $0 \le u < \sqrt{\alpha}$, we have $u^2 < \alpha \implies \alpha - u^2 > 0 \implies h'(u) > 0$. Therefore, the function is strictly increasing.
        \item For $u > \sqrt{\alpha}$, we have $u^2 > \alpha \implies \alpha - u^2 < 0 \implies h'(u) < 0$. Therefore, the function is strictly decreasing.
    \end{itemize}
    Since the function is strictly increasing before $u = \sqrt{\alpha}$ and strictly decreasing after, the critical point $u = \sqrt{\alpha}$ is a unique, global maximum of $h(u)$ on $[0, \infty)$.
    We evaluate $h(u)$ at this global maximum point:
    \begin{equation}
    h(\sqrt{\alpha}) = \frac{\sqrt{\alpha}}{(\sqrt{\alpha})^2 + \alpha} = \frac{\sqrt{\alpha}}{2\alpha} = \frac{1}{2\sqrt{\alpha}}
    \end{equation}
\end{itemize}
Consequently, we conclude that for all $x \ge 0$ and $y \ge 0$:
\begin{equation}
\gamma_c \le h(\sqrt{\alpha}) = \frac{1}{2\sqrt{\alpha}}
\end{equation}
This completes the step-by-step mathematical proof. 

Under our experimental choice of $\alpha = 0.05$, the spatial channel-wise scaling factor is analytically bounded by:
\begin{equation}
\gamma_c \le \frac{1}{2\sqrt{0.05}} \approx 2.236
\end{equation}
This bounds the scaling factor, mathematically neutralizing the Sparsity Trap by preventing scaling factor divergence and ensuring stable representation recovery across all data budgets.

\section{Experimental Details and Hyperparameters}
To guarantee full experimental reproducibility, we report our training and optimization hyperparameter details. All expert models are trained in a single GPU partition.

\subsection{Expert Fine-Tuning}
Each expert model (MNIST, Fashion-MNIST, and CIFAR-10) is initialized with ImageNet pre-trained weights. The backbone is a standard ResNet-18. We replace the original fully-connected layer with a newly initialized classification head consisting of a Dropout layer with rate $0.1$ followed by a Linear layer with 10 output classes.
Fine-tuning is performed for $5$ epochs using the standard SGD optimizer with momentum $0.9$ and weight decay $10^{-4}$. We employ a constant learning rate of $10^{-4}$ and a batch size of $128$.

\subsection{Evaluation Protocols}
All test evaluation accuracy values are computed using the full 10,000-sample test sets of each dataset. We pre-load the datasets onto the GPU memory once at startup to optimize computation speed. For the calibration sweeps, we evaluate across 5 distinct random seeds ($42, 43, 44, 45, 46$) to report robust statistics.

\bibliography{example_paper}
\bibliographystyle{icml2026}

\end{document}
